\documentclass[otchet]{../SCWorks}
% Тип обучения (одно из значений):
%    bachelor   - бакалавриат (по умолчанию)
%    spec       - специальность
%    master     - магистратура
% Форма обучения (одно из значений):
%    och        - очное (по умолчанию)
%    zaoch      - заочное
% Тип работы (одно из значений):
%    coursework - курсовая работа (по умолчанию)
%    referat    - реферат
%  * otchet     - универсальный отчет
%  * nirjournal - журнал НИР
%  * digital    - итоговая работа для цифровой кафедры
%    diploma    - дипломная работа
%    pract      - отчет о научно-исследовательской работе
%    autoref    - автореферат выпускной работы
%    assignment - задание на выпускную квалификационную работу
%    review     - отзыв руководителя
%    critique   - рецензия на выпускную работу

% * Добавлены вручную. За вопросами к @mchernigin
\usepackage{../preamble}

\begin{document}

% Кафедра (в родительном падеже)
\chair{математической кибернетики и компьютерных наук}

% Тема работы
\title{Анализ сложности для красно"=чёрного дерева}

% Курс
\course{2}

% Группа
\group{251}

% Факультет (в родительном падеже) (по умолчанию "факультета КНиИТ")
\department{факультета компьютерных наук и информационных технологий}

% Специальность/направление код - наименование
% \napravlenie{02.03.02 "--- Фундаментальная информатика и информационные технологии}
% \napravlenie{02.03.01 "--- Математическое обеспечение и администрирование информационных систем}
% \napravlenie{09.03.01 "--- Информатика и вычислительная техника}
\napravlenie{09.03.04 "--- Программная инженерия}
% \napravlenie{10.05.01 "--- Компьютерная безопасность}

% Для студентки. Для работы студента следующая команда не нужна.
\studenttitle{студентки}

% Фамилия, имя, отчество в родительном падеже
\author{Потапкиной Маргариты Андреевны}

% Заведующий кафедрой 
\chtitle{доцент, к.\,ф.-м.\,н.}
\chname{С.\,В.\,Миронов}

% Руководитель ДПП ПП для цифровой кафедры (перекрывает заведующего кафедры)
% \chpretitle{
%     заведующий кафедрой математических основ информатики и олимпиадного\\
%     программирования на базе МАОУ <<Ф"=Т лицей №1>>
% }
% \chtitle{г. Саратов, к.\,ф.-м.\,н., доцент}
% \chname{Кондратова\, Ю.\,Н.}

% Научный руководитель (для реферата преподаватель проверяющий работу)
\satitle{старший преподаватель} %должность, степень, звание
\saname{М.\,И.\,Сафрончик}

% Руководитель практики от организации (руководитель для цифровой кафедры)
\patitle{доцент, к.\,ф.-м.\,н.}
\paname{И.\,А.\,Батраева}

% Руководитель НИР
\nirtitle{доцент, к.\,п.\,н.} % степень, звание
\nirname{И.\,А.\,Батраева}

% Семестр (только для практики, для остальных типов работ не используется)
\term{2}

% Наименование практики (только для практики, для остальных типов работ не
% используется)
\practtype{учебная}

% Продолжительность практики (количество недель) (только для практики, для
% остальных типов работ не используется)
\duration{2}

% Даты начала и окончания практики (только для практики, для остальных типов
% работ не используется)
\practStart{01.07.2024}
\practFinish{13.01.2024}

% Год выполнения отчета
\date{2026}

\maketitle

% Включение нумерации рисунков, формул и таблиц по разделам (по умолчанию -
% нумерация сквозная) (допускается оба вида нумерации)
\secNumbering

\tableofcontents

% Раздел "Обозначения и сокращения". Может отсутствовать в работе
% \abbreviations
% \begin{description}
%     \item ... "--- ...
%     \item ... "--- ...
% \end{description}

% Раздел "Определения". Может отсутствовать в работе
% \definitions

% Раздел "Определения, обозначения и сокращения". Может отсутствовать в работе.
% Если присутствует, то заменяет собой разделы "Обозначения и сокращения" и
% "Определения"
% \defabbr

\intro
В данном отчёте производится анализ сложности и расхода памяти, а также описание особенностей красно"=чёрного дерева.

\section{Анализ сложности операций и расхода памяти}

\small
\inputminted[breaklines]{cpp}{./rb_tree.cpp}
\normalsize

Красно"=чёрное дерево поддерживает стандартные операции, присущие двоичному дереву поиска: создание нового узла, вставка и удаление узла, поиск минимального и максимального элемента, поиск элемента по значению, а также обходы: прямой, обратный и симметричный. Помимо этого, оно реализует некоторые специфичные функции, такие как левый и правый поворот, замена, получение указателя на брата, деда и дядю "--- эти функции применяются для перебалансировки дерева.

Ниже приводится анализ асимптотики этих операций.

\begin{itemize}
	\item создание узла, в т.ч. корня "--- $O(1)$;
	\item левый и правый повороты "--- $O(1)$;
	\item получение деда, дяди и брата "--- $O(1)$;
	\item вставка элемента "--- $O(h)$;
	\item 5 функций для разбора частных случаев при вставке элемента "--- $O(1)$;
	\item удаление элемента "--- $O(h)$;
	\item 6 функций для разбора частных случаев при удалении элемента "--- $O(1)$;
	\item замена "--- $O(1)$;
	\item поиск максимума и минимума "--- $O(h)$;
	\item поиск элемента по значению "--- $O(h)$;
	\item симметричный, прямой и обратный обход "--- $O(n)$;
\end{itemize}

Обходы работают за время $O(n)$, поскольку в них мы обрабатываем все узлы дерева, проходясь по каждому один раз. Создание узла требует $O(1)$ времени, при этом расход памяти на создание дерева составляет $O(5n) = O(n)$: в каждом узле хранится число, цвет узла (чёрный или красный), указатель на родителя и указатели на левого и правого сыновей. Остальные же функции выполняются за $O(h)$, поскольку в худшем случае мы вынуждены спускаться до листа.

Здесь:
\begin{itemize}
	\item[$\bullet$] $n$ "--- количество элементов в дереве;
	\item[$\bullet$] $h$ "--- высота бинарного дерева.
\end{itemize}

\section{Описание особенностей}

Красно"=чёрное дерево является частным случаем двоичного дерева поиска. Ключевая его особенность заключается в самобалансировке. В каждый момент времени красно"=чёрное дерево должно обладать следующими свойствами.

\begin{enumerate}
	\item Корень всегда чёрный.
	\item У красного узла чёрный родитель.
	\item Все простые пути из любого узла до листьев содержат одинаковое количество чёрных узлов. Это количество назовём <<чёрной высотой>>.
\end{enumerate}

В красно"=чёрном дереве с чёрной высотой $x$ не менее $2^{x - 1} - 1$ вершин. Отсюда следует, что дерево имеет логарифмическую высоту. В результате операции вставки, поиска и удаления работают за $O(log(n))$.

Соблюдение этих свойств достигается с помощью перекраски вершин и выполнения левого или правого поворота, когда происходит вставка вершины в дерево или её удаление.

% \conclusion

% Отобразить все источники. Даже те, на которые нет ссылок.
% \nocite{*}

\bibliographystyle{ugost2008}
\bibliography{thesis}

% Окончание основного документа и начало приложений Каждая последующая секция
% документа будет являться приложением
\appendix

\end{document}
